% Options for packages loaded elsewhere
\PassOptionsToPackage{unicode}{hyperref}
\PassOptionsToPackage{hyphens}{url}
%
\documentclass[
]{article}
\usepackage{amsmath,amssymb}
\usepackage{iftex}
\ifPDFTeX
  \usepackage[T1]{fontenc}
  \usepackage[utf8]{inputenc}
  \usepackage{textcomp} % provide euro and other symbols
\else % if luatex or xetex
  \usepackage{unicode-math} % this also loads fontspec
  \defaultfontfeatures{Scale=MatchLowercase}
  \defaultfontfeatures[\rmfamily]{Ligatures=TeX,Scale=1}
\fi
\usepackage{lmodern}
\ifPDFTeX\else
  % xetex/luatex font selection
\fi
% Use upquote if available, for straight quotes in verbatim environments
\IfFileExists{upquote.sty}{\usepackage{upquote}}{}
\IfFileExists{microtype.sty}{% use microtype if available
  \usepackage[]{microtype}
  \UseMicrotypeSet[protrusion]{basicmath} % disable protrusion for tt fonts
}{}
\makeatletter
\@ifundefined{KOMAClassName}{% if non-KOMA class
  \IfFileExists{parskip.sty}{%
    \usepackage{parskip}
  }{% else
    \setlength{\parindent}{0pt}
    \setlength{\parskip}{6pt plus 2pt minus 1pt}}
}{% if KOMA class
  \KOMAoptions{parskip=half}}
\makeatother
\usepackage{xcolor}
\usepackage{longtable,booktabs,array}
\usepackage{calc} % for calculating minipage widths
% Correct order of tables after \paragraph or \subparagraph
\usepackage{etoolbox}
\makeatletter
\patchcmd\longtable{\par}{\if@noskipsec\mbox{}\fi\par}{}{}
\makeatother
% Allow footnotes in longtable head/foot
\IfFileExists{footnotehyper.sty}{\usepackage{footnotehyper}}{\usepackage{footnote}}
\makesavenoteenv{longtable}
\setlength{\emergencystretch}{3em} % prevent overfull lines
\providecommand{\tightlist}{%
  \setlength{\itemsep}{0pt}\setlength{\parskip}{0pt}}
\setcounter{secnumdepth}{-\maxdimen} % remove section numbering
\ifLuaTeX
  \usepackage{selnolig}  % disable illegal ligatures
\fi
\IfFileExists{bookmark.sty}{\usepackage{bookmark}}{\usepackage{hyperref}}
\IfFileExists{xurl.sty}{\usepackage{xurl}}{} % add URL line breaks if available
\urlstyle{same}
\hypersetup{
  pdftitle={Word Embeddings},
  hidelinks,
  pdfcreator={LaTeX via pandoc}}

\title{Word Embeddings}
\author{}
\date{}

\begin{document}
\maketitle

Word embeddings are \textbf{dense vectors} that capture semantic and
syntactic word relationships by mapping words to a continuous vector
space. Unlike sparse representations (e.g., BoW or TF-IDF), embeddings
are trained to place similar words closer together in this space.

\hypertarget{key-properties}{%
\subsubsection{\texorpdfstring{\textbf{Key
Properties}}{Key Properties}}\label{key-properties}}

\begin{enumerate}
\tightlist
\item
  \textbf{Dimensionality Reduction}:

  \begin{itemize}
  \tightlist
  \item
    Compresses words into dense vectors (typically 50--300 dimensions
    vs. thousands in BoW).
  \end{itemize}
\item
  \textbf{Distributional Learning}:

  \begin{itemize}
  \tightlist
  \item
    Follows the \emph{distributional hypothesis} (like matrix methods)
    but uses neural networks to encode patterns.
  \end{itemize}
\item
  \textbf{Semantic Regularities}:

  \begin{itemize}
  \tightlist
  \item
    Captures analogies (e.g., \emph{king -- man + woman ≈ queen}) and
    hierarchies (e.g., \emph{dog → animal}).
  \end{itemize}
\end{enumerate}

\begin{longtable}[]{@{}lll@{}}
\toprule\noalign{}
Method & Pros & Limitations \\
\midrule\noalign{}
\endhead
\bottomrule\noalign{}
\endlastfoot
\textbf{BoW/TF-IDF} & Simple, interpretable & Sparse, no word
relationships \\
\textbf{LSA/PPMI} & Captures global stats & Linear, ignores word
order \\
\textbf{Embeddings} & Dense, nonlinear, context-aware & Requires large
data, less interpretable \\
\end{longtable}

\hypertarget{neural-vs.-count-based-approaches}{%
\subsubsection{\texorpdfstring{\textbf{Neural vs. Count-Based
Approaches}}{Neural vs. Count-Based Approaches}}\label{neural-vs.-count-based-approaches}}

While \textbf{matrix factorisation} (e.g., LSA, PPMI) relies on
co-occurrence statistics, neural embeddings (e.g., Word2Vec, GloVe)
\textbf{predict words from contexts} (or vice versa), enabling:

\begin{itemize}
\tightlist
\item
  \textbf{Nonlinear relationships} (e.g., polysemy handling).
\item
  \textbf{Efficient similarity computation} (cosine distance).
\item
  \textbf{Transfer learning} (pre-trained embeddings for downstream
  tasks).
\end{itemize}

\hypertarget{types-of-context}{%
\subsubsection{\texorpdfstring{\textbf{Types of
Context}}{Types of Context}}\label{types-of-context}}

Embeddings can be trained using:

\begin{itemize}
\tightlist
\item
  \textbf{Window-based contexts} (Word2Vec: local neighbours).
\item
  \textbf{Global co-occurrence} (GloVe: whole-corpus statistics).
\item
  \textbf{Sub-word information} (FastText: character n-grams).
\end{itemize}

\begin{center}\rule{0.5\linewidth}{0.5pt}\end{center}

\hypertarget{evaluation}{%
\section{Evaluation}\label{evaluation}}

Intrinsic (Tomas Mikolov, Kai Chen, Greg Corrado, Jeffrey Dean,
"Efficient Estimation of Word Representations in Vector Space")

\begin{itemize}
\tightlist
\item
  8869 semantic and 10675 syntactic questions -- cosine similarity
  between words
\end{itemize}

\textbf{1. Semantic-Syntactic Word Relationship Test Set}

\begin{itemize}
\tightlist
\item
  \textbf{Task Design}:

  \begin{itemize}
  \tightlist
  \item
    Created a comprehensive test set with \textbf{8,869 semantic} and
    \textbf{10,675 syntactic} questions.
  \item
    Examples:

    \begin{itemize}
    \tightlist
    \item
      \emph{Semantic}: "Athens is to Greece as Oslo is to Norway"
      (capital cities).
    \item
      \emph{Syntactic}: "big is to bigger as small is to smaller"
      (morphological patterns).
    \end{itemize}
  \end{itemize}
\item
  \textbf{Metric}:\\
  - \textbf{Accuracy}: A question is correct only if the predicted word
  (via vector arithmetic, e.g.,
  vec("King")−vec("Man")+vec("Woman")≈vec("Queen")vec("King")−vec("Man")+vec("Woman")≈vec("Queen"))
  exactly matches the ground truth. Synonyms are counted as errors.\\
  \textbf{2. Word Similarity Tasks}
\item
  \textbf{Datasets}: Evaluated on standard benchmarks (e.g., WordSim353,
  MEN) using \textbf{cosine similarity} between word vectors.
\item
  \textbf{Metric}: \textbf{Spearman's rank correlation} between model
  rankings and human-judged similarity scores.\\
  \textbf{3. Model Comparison}
\item
  \textbf{Baselines}: Compared against:

  \begin{itemize}
  \tightlist
  \item
    Traditional NNLM (Feed-forward and RNN-based).
  \item
    Publicly available word vectors (e.g., Collobert-Weston, Turian).
  \end{itemize}
\item
  \textbf{Key Results}:

  \begin{itemize}
  \tightlist
  \item
    Skip-gram outperformed CBOW on \textbf{semantic tasks} (55\% vs.
    24\% accuracy).
  \item
    CBOW was better for \textbf{syntactic tasks} (64\% vs. 59\%).
  \item
    Both models surpassed NNLM and RNNLM in accuracy and training speed
    (Table 3).\\
    \textbf{4. Scaling Experiments}
  \end{itemize}
\item
  \textbf{Training Data Impact}: Tested on subsets of Google News corpus
  (24M to 6B words).

  \begin{itemize}
  \tightlist
  \item
    Findings: Larger data + higher dimensionality (d=300−1000d=300−1000)
    improved accuracy, but with diminishing returns (Table 2).
  \end{itemize}
\item
  \textbf{Efficiency}:\\
  - CBOW trained faster (1 day for 6B words) vs. Skip-gram (3 days).\\
  - Distributed training (DistBelief) scaled to \textbf{trillions of
  words} (Table 6).\\
  \textbf{5. Microsoft Sentence Completion Challenge}
\item
  \textbf{Task}: Predict missing words in sentences from 5 choices.
\item
  \textbf{Result}: Skip-gram alone achieved \textbf{48\% accuracy};
  combined with RNNLM, reached \textbf{58.9\%} (new SOTA, Table 7).\\
  \textbf{6. Qualitative Analysis}
\item
  \textbf{Vector Arithmetic}: Demonstrated relationships like:

  \begin{itemize}
  \tightlist
  \item
    vec("Paris")−vec("France")+vec("Italy")≈vec("Rome")vec("Paris")−vec("France")+vec("Italy")≈vec("Rome").
  \item
    Analogies for grammar (e.g., verb tenses) and semantics (e.g.,
    currencies).
  \end{itemize}
\end{itemize}

Qualitative evaluation: \url{https://projector.tensorflow.org/}\\
Extrinsic

\begin{itemize}
\tightlist
\item
  Ex. compare word embeddings and n-grams as input to a number of
  downstream\\
  tasks
\end{itemize}

\hypertarget{visualising-word-embeddings}{%
\section{Visualising Word
Embeddings}\label{visualising-word-embeddings}}

\begin{itemize}
\tightlist
\item
  Word embeddings can be used to visualise the meaning of a word by:
\item
  listing the words in V with highest cosine similarity with the word.
\item
  project the d dimensions of a word embedding down into 2 dimensions
  (t-distributed stochastic neighbor embedding (t-SNE)).
\end{itemize}

\hypertarget{properties-and-types}{%
\section{Properties and Types}\label{properties-and-types}}

Ability to capture relational meanings.\\
Can be used to solve anlogy provlems with a parallelogram model:\\
\url{https://doi.org/10.1016/j.cognition.2020.104440}

\hypertarget{diachronic-word-emebeddings}{%
\paragraph{Diachronic Word
Emebeddings}\label{diachronic-word-emebeddings}}

\url{https://en.wikipedia.org/wiki/Orthogonal_Procrustes_problem}\strut \\
Comparing embeddings trained on different but related data.\\
same domain different time periods.\\
Orthogonal Procrustes Problem: how to align embeddings into same space
to be able to do meaningful comparisons.

{}Where, Q is a transformation (rotation and alignment) matrix.\\
Solution is {} where {} are from SVD decomposition of {}.

\hypertarget{polysemous-word-embeddings}{%
\paragraph{Polysemous Word
Embeddings}\label{polysemous-word-embeddings}}

Embeddings of words with mulitple meaings;\\
Arora et al. ``Linear Algebraic Structure of Word Senses, with
Applications to Polysemy''
(\url{https://www.aclweb.org/anthology/Q18-1034.pdf})

Shows that the resulting embeddings will be proportional to the sum of
the embeddings of the various meanings.\\
v(tie) =\textasciitilde{} a1v(tie1) + a2v(tie2) + a3v(tie3)\\
where ai is proportional to the relative frequency of the meaning i\\
• v(tie1) is a discourse atom (\textasciitilde topic)\\
that can be learned (here represented\\
by a set of representative words)

Properties\\
\textbf{Methods} (Word2Vec, GloVe, FastText, CBOW)

\hypertarget{word2vec}{%
\paragraph{Word2Vec}\label{word2vec}}

Major acheivement:

\begin{itemize}
\tightlist
\item
  defining a learning problem without the need to annotate data manually
  (overcomes major limiting factor for machine learning)
\item
  unlocks massive datasets
\end{itemize}

Train a classifier on a binary prediction task:

\begin{itemize}
\tightlist
\item
  Given a word {} predict if a word {} is in the context\\
  - generating term-context {} pairs\\
  Thus we want our classifier to estimate the following
  probability:{}Whereby {}{}Assuming independence of context words we
  can compute the log probability\\
  {}
\end{itemize}

Learning function that maximising probability of context words and
minimises non-context words\\
\${}

Training process: 1. treat target word and neighbouring context as
positive examples 2. randomly sample other words in lexicon for negative
samples 3. logistic regression to train classifier 4. use the learned
weights as the embeddings Model perspective: - each word t in corpus has
two d-dimensional representations totalling \$d*2\textbar V\textbar\$
parameters - w - target words representation - c - context and noise
word representation - final embedding form for word i: 1) \$w\_i +
c\_i\$ 2) \$w\_i\$ Gradient of loss function: \$\$w\^{}\{t+1\} = w\^{}t
- \textbackslash eta
{[}\textbackslash sigma(c\_\{pos\}\textbackslash cdot
w\^{}t)-1{]}c\_\{pos\} +
\textbackslash sum\_\{i=1\}\^{}k{[}\textbackslash sigma(c\_\{neg\_i\}\textbackslash cdot
w\^{}t){]}c\_\{neg\_i\}{]}\$\$ - \$w\^{}\{t+1\}\$, \$w\^{}t\$: next word
to predict t+1 using current word t. - **η**: learning rate - **The
bracketed term**: The gradient of the loss function with respect to the
word vector. 1. positive pair (word t and neighbouring context):
\$\textbackslash sigma(c\_\{pos\}\textbackslash cdot
w\^{}t)-1{]}c\_\{pos\}\$ - cpos\hspace{0pt}: This is the vector
representation of the positive context word. - wt⋅cpos\hspace{0pt}: This
is the dot product between the target word vector and the positive
context word vector. A higher dot product indicates greater similarity.
- (σ(cpos\hspace{0pt}⋅wt)−1): Sigmoid to squash dot between 0 and 1.
Ideally, for a true positive pair (which it should be since its real
context), we want σ(cpos\hspace{0pt}⋅wt) to be close to 1. If
it\textquotesingle s less than 1, this term will be negative, pushing
the word vectors to become more similar. - The entire term
(σ(cpos\hspace{0pt}⋅wt)−1)cpos\hspace{0pt} contributes to adjusting wt
in the direction that increases its similarity with the positive context
word cpos\hspace{0pt}. 2.
**∑i=1k\hspace{0pt}σ(cnegi\hspace{0pt}\hspace{0pt}⋅wt)cnegi\hspace{0pt}\hspace{0pt}**:
This part relates to the \_negative samples\_. - k: num negative
hyperparamater samples - cnegi\hspace{0pt}\hspace{0pt}: This is the
vector representation of the i-th negative context word (a word that is
\_not\_ actually in the context of w). -
wt⋅cnegi\hspace{0pt}\hspace{0pt}: This is the dot product between the
target word vector and the i-th negative context word vector. -
σ(cnegi\hspace{0pt}\hspace{0pt}⋅wt): This is the sigmoid output,
representing the probability that cnegi\hspace{0pt}\hspace{0pt} is a
context word for w. For negative samples, we want this probability to be
close to 0. - The entire summation term contributes to adjusting wt in
the direction that decreases its similarity with the negative context
words cnegi\hspace{0pt}\hspace{0pt}. **More on Negative sampling:** For
each positive pair (target word, context word), we create \$k\$
(hyperparameter) negative pairs by randomly selecting a context word
with probability \$P\_\textbackslash alpha(w)\$ (where
\$\textbackslash alpha\$ is a hyperparameter, usually set to 0.75,
increasing the probability of rare words). - The probability of
selecting a word \$w\$ as a negative sample is given by:
\$\$P\_\textbackslash alpha(w) =
\textbackslash frac\{count(w)\^{}\textbackslash alpha\}\{\textbackslash sum\_\{w\textquotesingle{}
\textbackslash in Vocabulary\}
count(w\textquotesingle)\^{}\textbackslash alpha\}\$\$where -
\$count(w)\$ is the frequency of word \$w\$ in the training corpus. -
\$\textbackslash alpha\$ is a hyperparameter (typically 0.75). - The
exponent \$\textbackslash alpha \textless{} 1\$ increases the
probability of selecting rare words. - Reduces the computational cost by
considering only a small number of negative examples instead of the
entire vocabulary. **Obersations:** - Context words randomly sampled
rfavouirng closer words with hyperparamter window size L - L changes the
utility of our embeddings: - Small window size = semantically similar
with the same positions - larger window size = topical relatedness -
Pre-trained word embeddings are the most popular represenation for
learning algoriths - be careful: bias may propagate throughout the
system **Connection to matrix factorizaiton:** Skip-Gram with negative
sampling model is connected to matrix factorization of C{[}1{]} : - Ci,j
= PMI(i,j) -- log k - PMI of word i in context j - k is the number of
negative examples sampled - We can use PPMI to avoid --infinite values
when count(i, j)=0
https://transacl.org/ojs/index.php/tacl/article/view/570/124 **Title:**
Improving Distributional Similarity with Lessons Learned from Word
Embeddings **Authors:** Omer Levy, Yoav Goldberg, Ido Dagan **Published
in:** Transactions of the Association for Computational Linguistics
(TACL), 2015 This work reshaped understanding of word embeddings by
highlighting the importance of hyperparameters and enabling fairer
comparisons between methods. \#\#\#\#\# **Key Contributions:** - The
paper reveals that much of the performance gains attributed to neural
word embeddings (e.g., Word2Vec, GloVe) stem from system design choices
and hyperparameter optimisations, rather than the underlying algorithms.
- The authors demonstrate that these optimizations (e.g., dynamic
context windows, subsampling, negative sampling, and context
distribution smoothing) can be applied to traditional count-based
methods like PPMI and SVD, closing the performance gap between the two
approaches. - When hyperparameters are properly tuned, count-based
methods (PPMI, SVD) perform comparably to neural embeddings (SGNS,
GloVe) on word similarity and analogy tasks. - Context distribution
smoothing (a technique borrowed from Word2Vec) significantly improves
PPMI by mitigating its bias toward rare word co-occurrences. - There is
no consistent global advantage of neural embeddings over count-based
methods when both are optimized. - **Practical Recommendations**: -
Always use context distribution smoothing (cds=0.75) for PPMI/SVD. -
Avoid the traditional SVD setup (eig=1); symmetric variants (eig=0 or
0.5) perform better. - SGNS is a robust baseline due to its efficiency
and competitive performance. \#\#\#\#\# **Contradictions to Prior
Work**: - Challenges claims that neural embeddings universally
outperform count-based methods (e.g., Baroni et al., 2014), showing that
differences often arise from untuned hyperparameters in traditional
methods. - Finds that GloVe does not outperform SGNS when both are
fairly compared, contrary to Pennington et al. (2014). \#\#\#\#
Fast-Text https://arxiv.org/pdf/1607.04606.pdf (https://fasttext.cc) A
solution for learning word embeddings for languages with a rich
morphological structure. This is done by computing representation for
char n-grams + \textless{} for word beggning and \textgreater{} for word
endings. Words are then the sum of the embeddings of the char n-grams
comprosig it. n is a hyperparamater that is dependant on the language
and the task. it also shows general improvements for syntatic analogy
tasks on rare words. \#\#\#\# Glove *Pennigton et al. ``Glove: Global
Vectors for Word Representation'' (2014)
https://aclanthology.org/D14-1162.pdf Count-based model for learning
word embeddings. Unlike the predictive methods like Skip-Gram and CBOW
in Word2Vec, which learn embeddings by predicting context words or a
target word, GloVe leverages the global co-occurrence statistics of
words within a corpus. Initialises embeddings based on the global
co-occurence statistics between words from the term-term matrix. **Ratio
of Co-occurrence Probabilities** The fundamental intuition behind GloVe
is that the ratios of word-word co-occurrence probabilities have the
potential to encode meaning. For example, consider two target words,
"ice" and "steam," and two context words, "solid" and "gas." - "ice" is
expected to co-occur frequently with "solid" and infrequently with
"gas." - "steam" is expected to co-occur frequently with "gas" and
infrequently with "solid." - The ratio of co-occurrence probabilities
\$\textbackslash frac\{P(solid \textbar{} ice)\}\{P(gas \textbar{}
ice)\}\$ should be large. - The ratio of co-occurrence probabilities
\$\textbackslash frac\{P(solid \textbar{} steam)\}\{P(gas \textbar{}
steam)\}\$ should be small. - The ratio \$\textbackslash frac\{P(solid
\textbar{} ice)\}\{P(solid \textbar{} steam)\}\$ should be large. - The
ratio \$\textbackslash frac\{P(gas \textbar{} steam)\}\{P(gas \textbar{}
ice)\}\$ should be large. GloVe aims to learn word vectors such that
their dot product is related to the logarithm of these co-occurrence
probabilities. **Methodology:** 1. **Constructing the Co-occurrence
Matrix (X):** - GloVe first scans the entire corpus to build a large
word-word co-occurrence matrix \$X\$, where \$X\_\{ij\}\$ represents the
number of times word \$j\$ appears in the context of word \$i\$ - The
context is typically defined by a symmetric window around the target
word. - A weighting function \$f(X\_\{ij\})\$ is often used to
down-weight the contribution of very frequent co-occurrences, as these
might be less informative. A common form for \$f(x)\$ is: \$\$f(x) =
\textbackslash begin\{cases\} (x/x\_\{max\})\^{}\textbackslash alpha \&
\textbackslash text\{if \} x \textless{} x\_\{max\}
\textbackslash\textbackslash{} 1 \& \textbackslash text\{otherwise\}
\textbackslash end\{cases\}

\begin{verbatim}
    where $x_{max}$ and $\alpha$ are hyperparameters.
\end{verbatim}

\begin{enumerate}
\setcounter{enumi}{1}
\item
  \textbf{Defining the Loss Function:}

  \begin{itemize}
  \tightlist
  \item
    Learn two sets of word vectors: {} (word vector of word {}) and {}
    (context word vector of word {}).
  \item
    Minimise the following weighted least squares loss function:\\
    {} where:

    \begin{itemize}
    \tightlist
    \item
      {} is the size of the vocabulary.
    \item
      {} and {} are bias terms associated with words {} and {},
      respectively.
    \item
      The term {} comes from the idea of matching the dot product of
      word vectors to the logarithm of the co-occurrence counts.
    \end{itemize}
  \end{itemize}
\item
  Train with SGD.
\item
  \textbf{Final Embeddings:}

  \begin{itemize}
  \tightlist
  \item
    same as Word2Vec
  \item
    After training, the final word embedding for a word {} is typically
    obtained by summing its word vector and its context word vector: {},
    or simply using {}
  \end{itemize}
\end{enumerate}

\textbf{Advantages of GloVe:}

\begin{itemize}
\tightlist
\item
  \textbf{Leverages Global Statistics:} Directly using the co-occurrence
  matrix global relationships between words are captured, local context
  window methods might miss broader semantic connections.
\item
  \textbf{Strong Empirical Performance:} Often competitive with or
  superior to Word2Vec models.
\end{itemize}

\textbf{Disadvantages:}

\begin{itemize}
\tightlist
\item
  \textbf{Computational Cost of Co-occurrence Matrix:} Building the
  co-occurrence matrix can be computationally expensive for very large
  corpora.
\item
  \textbf{Less Flexibility with Context:} Unlike Word2Vec, where the
  context window can be adjusted during training, the co-occurrence
  statistics in GloVe are pre-computed based on a fixed window size.
\end{itemize}

\hypertarget{continuous-bag-of-words-model-cbow}{%
\paragraph{Continuous Bag of Words model
(CBOW)}\label{continuous-bag-of-words-model-cbow}}

The Continuous Bag of Words (CBOW) model employs a sliding window
training technique. The window consists of a central target word and its
surrounding context words (both forward and backward).

At each training step, the model predicts the probability distribution
of the centre word given its context. The model parameters are then
updated using Stochastic Gradient Descent (SGD).

CBOW is typically implemented as a fully connected feed-forward neural
network with the following key hyperparameters:

\begin{itemize}
\tightlist
\item
  {}: Word embedding size. This dimension also implicitly influences the
  effective context window size considered during training.
\end{itemize}

The network architecture involves:

\begin{itemize}
\tightlist
\item
  An input layer representing the context words.
\item
  One or more hidden layers with a non-linear activation function,
  commonly Rectified Linear Unit (ReLU):\\
  {}
\item
  An output layer that produces a probability distribution over the
  entire vocabulary using the softmax function:\\
  {}where {} is the output vector from the previous layer and {} is the
  vocabulary size.
\end{itemize}

The model is trained by minimizing the cross-entropy loss function,
which measures the dissimilarity between the predicted probability
distribution and the one-hot encoded representation of the actual center
word. The cross-entropy loss {} for a single training example is given
by:\\
{}where {} is the one-hot encoded target word and {} is the predicted
probability distribution.

The input to the model consists of the context words. The center word is
represented as a one-hot encoded vector. The context word vector used as
input to the neural network is typically the average of the one-hot
encoded vectors of all the context words within the sliding window.

Thus, \textbf{word order within the context is not explicitly modelled},
similar to the Bag-of-Words (BoW) approach.

The weight matrices of the neural network learn the word embeddings:

\begin{enumerate}
\tightlist
\item
  Each column of the first weight matrix, {}, represents the input
  embedding vector of the corresponding word in the vocabulary (where {}
  is the vocabulary size).
\item
  Each row of the second weight matrix, {}, represents the output
  embedding vector of the corresponding word in the vocabulary.
\item
  The final word embedding for each word is often taken as the average
  of its input embedding (from {}) and the transpose of its output
  embedding (from {}).
\end{enumerate}

\hypertarget{other-methods}{%
\paragraph{Other Methods}\label{other-methods}}

\begin{itemize}
\item
  \textbf{WordNet}: A manually curated lexical database that groups
  words into sets of synonyms (synsets), provides short, general
  definitions (glosses), and records various semantic relations between
  these synsets.

  \begin{itemize}
  \tightlist
  \item
    \textbf{Limitation}: WordNet is often incomplete for most real-world
    applications due to the vastness and dynamic nature of language.
  \end{itemize}
\item
  \textbf{First-Order Logic (FOL)} (Eisenstein 12.2): A formal system of
  symbolic logic used in mathematics, philosophy, linguistics, and
  computer science. It extends propositional logic by allowing
  quantification over objects and the use of predicates to express
  properties and relations.

  \begin{itemize}
  \item
    \textbf{Proposition}: A declarative sentence that can be either true
    or false and is associated with a truth value.
  \item
    \textbf{Denotation of a proposition}: The set of all equivalent
    forms of a proposition that share the same truth value.

    \begin{itemize}
    \tightlist
    \item
      Example: {} (all denote the truth value
      \textquotesingle true\textquotesingle{} if interpreted as
      mathematical equalities).
    \end{itemize}
  \item
    \textbf{Boolean operators}: Logical connectives that combine
    propositions to form more complex propositions (e.g., and ({}), or
    ({}), not ({}), implication ({}), equivalence ({})).
  \item
    \textbf{Quantifiers}: Symbols that express the extent to which a
    predicate is true over a range of elements:

    \begin{itemize}
    \tightlist
    \item
      \textbf{Universal quantifier}: "For all" ({}).
    \item
      \textbf{Existential quantifier}: "There exists" ({}).
    \end{itemize}
  \item
    \textbf{Constants, variables, and relations}: The basic building
    blocks of FOL:

    \begin{itemize}
    \tightlist
    \item
      \textbf{Constants}: Symbols that represent specific objects (e.g.,
      Italy, Rome).
    \item
      \textbf{Variables}: Symbols that can represent any object in the
      domain of discourse (e.g., {}, {}).
    \item
      \textbf{Relations (Predicates)}: Symbols that express
      relationships between objects or properties of objects (e.g.,
      Capital({}, {}), is\_a({}, country)).
    \end{itemize}
  \item
    \textbf{Example of a relation represented as a set of tuples}:\\
    {} This relation represents the capital city of each country in the
    set.
  \item
    \textbf{Inference and discovery of new facts}: First-order logic
    provides a formal framework for making logical inferences based on
    existing knowledge, allowing for the derivation and discovery of new
    facts that are logically entailed by the initial set of axioms and
    rules.
  \end{itemize}
\end{itemize}

\textbf{Optimization} (Negative sampling, hierarchical softmax)\\
\textbf{Evaluation} (Intrinsic/extrinsic tasks)

\end{document}
